\section{Finite lengths, asymptotics, and complexity}
\label{sec:scaling}

Finite validity, asymptotic distance, and encoding complexity are separate
properties. We state them separately for each SPIN family.

\subsection{Native and requested lengths}

The direct constructions apply at native lengths for which their block
dimensions are integral. Their logarithmic schedules permit admissible lengths
with relative gaps \(o(1)\). For Structured SPIN, the zero extension in
\eqref{eq:structured-requested-length-wrapper} maps the largest native member
below any sufficiently large requested length into that output length. This wrapper preserves
injectivity and absolute minimum distance. It changes rate and relative
distance by only $o(1)$ and preserves linear work.

\subsection{Asymptotic statements}

Theorems~\ref{thm:independent-accumulator-spin}
and~\ref{thm:independent-random-spin} give positive relative distance with
high probability for the independent-block ensembles. They specify the block
schedule, setup failure probability, and rate. At rate one half, the stated
examples use block constants \(3\) and \(17\), respectively; Random SPIN also
uses memory constant \(51/50\).

Theorem~\ref{thm:structured-spin-scalable} uses direct lengths
$N_m=L_mb_m$, where $L_m=128m$ and $b_m$ is the least multiple of $24$
satisfying $b_m\ge\frac{39}{4}\log_2N_m$. These lengths have relative gaps $o(1)$.
The requested-length wrapper gives rate $1/2-o(1)$, relative distance
$0.11-o(1)$, and linear work.

\subsection{Finite parameter selection}
\label{sec:finite-parameter-extrapolation}

The finite theorem and Figure~\ref{fig:finite-bch-curve} provide certified
choices, rather than an extrapolated parameter schedule.
The selector's objective is measured encoding cost subject to a complete
first-moment bound at the requested length and margin.
A diagnostic calculation can rank candidates, but acceptance requires
outward arithmetic and coverage of every nonzero occupancy.

The inner parameters control different costs. Increasing $t$ reduces the
number of recursive steps; increasing $s$ can suppress cancellations but
increases the local arithmetic. The selected map also matters, not merely
its dimensions. For $t=2^m$, an $s$-dimensional subcode of
$\operatorname{RM}(2,m)$ requires
\begin{equation}
  s\le 1+m+\binom m2.
  \label{eq:rm2sub-capacity}
\end{equation}
This dimension constraint is an admissibility test, not a distance guarantee.
Changing the outer or selected map requires a matching spectrum calculation
and inner analysis. The fixed $(128,19)$ ladder shows that useful finite
coverage need not require retuning at every message length.

\subsection{Random-family encoding costs}

Under either proved schedule, the outer encoder processes \(N/B\) blocks of
length \(B=\Theta(\log N)\). Direct dense block multiplication therefore
costs \(\Theta(N\log N)\) bit operations. The accumulator costs \(O(N)\).
The Random SPIN recurrence evaluates \(m=\Theta(\log N)\) taps per output
coordinate and hence costs \(\Theta(N\log N)\). These bounds concern direct
encoders; preprocessing or word-level batching may reduce practical constants.

\subsection{Cost model}

Let $C_{\mathrm{out}}(b)$ be the work to encode one outer block, and let
$C_{\mathrm{in}}(t,s)$ be the work for one structured inner step. A
Structured SPIN encoder has the schematic cost
\begin{equation}
  T(N)
  =
  \frac{N}{b}C_{\mathrm{out}}(b)
  +
  \frac{N}{t}C_{\mathrm{in}}(t,s)
  +O(N),
  \label{eq:structured-cost}
\end{equation}
where the last term covers routing and permutation work.

For any fixed base constituent with a constant-size encoding circuit, tiling,
two permutations, and two accumulators give
$C_{\mathrm{out}}(b)=O(b)$.\linebreak The fixed RM2Sub-S19 step has constant cost.
Every term in \eqref{eq:structured-cost} is therefore $O(N)$. The current
certificate uses extended Golay as the base constituent, but the cost argument
does not depend on its 24-bit length.
